\chapter{Chord and arc}
\label{app:chord}

This appendix proves the inequality \eqref{eq:sinx},
\[
  \cos x<\frac{\sin x}{x}<1\qquad(0<x<\pi/2),
\]
used in \cref{ch:statespace} to compare the chord and the arc swept by a
turn error. It needs no derivative or series, only areas.

\begin{figure}[!htb]
\centering
\begin{tikzpicture}[scale=1.25,>=Stealth]
  \def\s{2.6}\def\x{50}
  \coordinate (P) at (0,0);
  \coordinate (A) at (0:\s);            % p'
  \coordinate (Q) at (\x:\s);           % q
  \coordinate (T) at ({\s},{\s*tan(\x)}); % t on the tangent at p'
  \fill[orange!30] (P) -- (A) -- (T) -- cycle;
  \fill[blue!25] (P) -- (A) arc[start angle=0,end angle=\x,radius=\s]
    -- cycle;
  \fill[green!40] (P) -- (A) -- (Q) -- cycle;
  \draw[black!30] (-8:\s) arc[start angle=-8,end angle=\x+12,radius=\s];
  \draw[thick] (P) -- (A) node[pos=0.5,below] {$s$};
  \draw[thick] (P) -- (T);
  \draw[thick,green!50!black] (A) -- (Q);
  \draw[densely dotted,black!60] ($(A)+(0,-0.5)$) -- ($(T)+(0,0.4)$)
    node[above,font=\small,black!60] {tangent at $\p'$};
  \draw[black!60] (A) ++(-0.2,0) -- ++(0,0.2) -- ++(0.2,0);
  \draw[->,black!70] (0.9,0) arc[start angle=0,end angle=\x,radius=0.9];
  \node[font=\small] at (\x/2:1.15) {$x$};
  \fill (P) circle (1.5pt) node[below left] {$\p$};
  \fill (A) circle (1.5pt) node[below right] {$\p'$};
  \fill (Q) circle (1.5pt) node[above left] {$\bm r$};
  \fill (T) circle (1.5pt) node[right] {$\bm t$};
  \draw[decorate,decoration={brace,amplitude=4pt,raise=2pt}]
    (A) -- (T) node[midway,right=7pt,font=\small] {$s\tan x$};
  \node[font=\small,anchor=west] at (4.1,2.0)
    {\tikz\fill[green!40] (0,0) rectangle (0.3,0.2); triangle $\p\,\p'\bm r$:
    $\tfrac12 s^2\sin x$};
  \node[font=\small,anchor=west] at (4.1,1.5)
    {\tikz\fill[blue!25] (0,0) rectangle (0.3,0.2); sector: $\tfrac12 s^2 x$};
  \node[font=\small,anchor=west] at (4.1,1.0)
    {\tikz\fill[orange!30] (0,0) rectangle (0.3,0.2); triangle $\p\,\p'\bm t$:
    $\tfrac12 s^2\tan x$};
\end{tikzpicture}
\caption{The three nested regions behind $\sin x<x<\tan x$. Each contains
the previous one, so their areas increase.}
\label{fig:sinx}
\end{figure}

Take a circle of radius $s$ centred at $\p$, a point $\p'$ on it, and the
point $\bm r$ on the circle seen from $\p$ at the angle $x$ from $\p'$
(\cref{fig:sinx}). Let $\bm t$ be the point where the ray from $\p$ through
$\bm r$ meets the tangent at $\p'$. Three regions are nested:
\begin{itemize}
  \item the triangle $\p\,\p'\,\bm r$, with base $\|\p-\bm r\|=s$ and height
    $s\sin x$ (the distance from $\p'$ to the line $\p\bm r$), so area
    $\tfrac12 s^2\sin x$;
  \item the circular sector $\p\,\p'\,\bm r$, a fraction $x/(2\pi)$ of the
    disc of area $\pi s^2$, so area $\tfrac12 s^2 x$;
  \item the triangle $\p\,\p'\,\bm t$, right-angled at $\p'$ with legs $s$
    and $s\tan x$, so area $\tfrac12 s^2\tan x$.
\end{itemize}
The first is inside the second because the sector is convex (its angle is
less than $\pi$) and contains the three vertices $\p,\p',\bm r$. The second
is inside the third because the whole disc lies on the same side of the
tangent as $\p$: any point $\bm y\ne\p'$ of the tangent satisfies
$\|\bm y-\p\|^2=s^2+\|\bm y-\p'\|^2>s^2$ (Pythagoras, the tangent being
perpendicular to the radius), so the tangent line never enters the disc.
Each inclusion leaves out a region of positive area, so
$\tfrac12 s^2\sin x<\tfrac12 s^2 x<\tfrac12 s^2\tan x$. Dividing by
$\tfrac12 s^2 x$ gives $\sin x/x<1$, and dividing by $\tfrac12 s^2\tan x$
gives $\cos x<\sin x/x$: this is \eqref{eq:sinx}. No derivative or series
is needed; the only analytic input is that the area of a sector is
proportional to its angle, which is how the radian is defined.
